%--------------------------------------------------------------------------------
%
%
%--------------------------------------------------------------------------------
\pagebreak
\section*{RUN — Run the simulation}
\addcontentsline{toc}{section}{RUN — Run the simulation}

%--------------------------------------------------------------------------------
\begin{iPageDescEntry}{Syntax}
\texttt{RUN} \\
\end{iPageDescEntry}

%--------------------------------------------------------------------------------
\begin{iPageDescEntry}{Description}

The \texttt{RUN} command starts continuous execution of the simulation. It is
functionally equivalent to repeatedly stepping the simulation without an
explicit step count.

Once invoked, control is transferred to the TWIN-64 system, and all configured
and enabled processors execute instructions continuously. Execution proceeds
until a condition occurs that returns control to the simulator.

Control is returned to the simulator when a trap, breakpoint exception, trace
exception, machine check, or other architectural exception is raised. At that
point, execution halts and the simulator command interface becomes active again,
allowing inspection or modification of system state.

Unlike the \texttt{STEP} command, \texttt{RUN} does not impose a limit on the
number of executed instructions and continues indefinitely until an exceptional
event occurs.

\end{iPageDescEntry}

%--------------------------------------------------------------------------------
\begin{iPageDescOpEntry}{Example}
\begin{lstlisting}[style=iPageOpStyle]

(80)->run          # start continuous simulation execution

\end{lstlisting}
\end{iPageDescOpEntry}

%--------------------------------------------------------------------------------
\begin{iPageDescEntry}{Notes}

The \texttt{RUN} command respects breakpoint and trace settings. If breakpoints or
single-step tracing are enabled, execution may stop shortly after \texttt{RUN}
is issued.

\end{iPageDescEntry}
